% Pursuing Stacks (cotes 134-2 to 134-8): archive cote 134-5, site page 29 (PDF page 30), typescript page 216.
% typescript. Draft compared: ps4.tex:1359-1429.
% AI-transcribed by Claude Opus 5.5 (claude-opus-5-5, adaptive thinking, effort high), Message Batches API, batch pilot (msgbatch_0189Nhf1xB615YWnztubtLfc), 2026-11-12; not yet checked by eye.
\scanpage{134-5/29}
\mspage{216}

(vi') Same as (vi), with F restricted to be an object b in B.

Indeed, we have of course (i)$\Rightarrow$(vi)$\Rightarrow$(vi'), but also (vi')$\Rightarrow$(iii)
if we admit strong saturation of $\underline{W}$, which implies that an object X of
(Cat) is $\underline{W}$-aspheric iff it's image in Hot($\underline{W}$) is a final object.

\uline{Remark} 3. If we don't assume A and B to be $\underline{W}$-aspheric, the implications
\[ \text{(As)} \Longrightarrow \text{(iv')} \Longrightarrow \text{(iv)} \]
are both strict. As an illustration of this point, take for B a discre-%
te category (defined in terms of a set of indices I = Ob B), thus a
\margin{$(A_b)_{b\in I}$} \xout{\illegible{2}} category A over B is essentially the same as a family $(A_{\hand{b}})_{\hand{b\in I}}$
of objects of (Cat) indexed by I. In terms of this\xout{e} family, we see at
once that the three conditions above on i: A$\rightarrow$B mean respectively
(a) that all categories $A_b$ are $\underline{W}$-aspheric, (b) that all categories
$A_b$ are non-empty, and (c) condition vacuous. This example brings to
mind that \uline{condition} (iv) \uline{is a Serre-fibration type condition} , we'll
come back upon this condition when studying \xout{\illegible{1}} homotopy properties
of maps in (Cat)\hand{,} with special emphasis on base change questions. Like-%
wise, condition (iv') appears as a strengthening of such Serre-type
condition, to the effect that the restriction of A over any connected
component of B should be moreover non-empty. As an example, we may take
the projection BxC $\rightarrow$ \over{\unsure{C}}{B} , where C is any non-empty object in (Cat).

I would like now to dwell still a little on the case, of special
interest of course for the modelizing story, when A and B are test cate-%
gories or something of the kind. The formulation (i) of the asphericity
condition \over{of}{for} the functor i can be expressed by stating that in the
the functor i\hand{$\overline{\phantom{i}}^{*}$} \xout{on} \deleted{the}\inserted{between} localizations $\mathrm{Hot}_B$ and $\mathrm{Hot}_A$ exists, and gives
rise (via \hand{$\lambda$}$_i$) to a commutation morphism which is an \uline{isomorphism}
\diagram{
\begin{tikzcd}[column sep=small]
\mathrm{Hot}_B \arrow[rr, "\hand{$\overline{i^*}$}"] \arrow[dr, swap, "\hand{$\overline{i_B}$}"] & & \mathrm{Hot}_A \arrow[dl, "\hand{$\overline{i_A}$}"] \\
& \mathrm{Hot}(\underline{W}) &
\end{tikzcd}
\quad \text{with, between the two slanted arrows, a handwritten arrow pointing left labelled } \hand{$\overline{\lambda_i}$}
}
\margin{$\overline{\lambda_i}$ over a leftward arrow marked $\sim$}
.

This shows, when A and B are pseudo-test categories for $\underline{W}$ , i.e. the
functors \hand{$\overline{i_A}$} and \hand{$\overline{i_B}$} are equivalences, that the functor i\hand{$\overline{\phantom{i}}^{*}$} deduced
from the \inserted{$\underline{W}$-}aspheric \xout{functor} map i: A$\rightarrow$B , \uline{does not depend} (up to cano-%
nical isomorphism) on the choice of i , and can be described as the com-%
position of \hand{$\overline{i_B}$} followed by a quasi-inverse for \hand{$\overline{i_A}$} . This of course
is very much in keeping with the ``inspiring assumption'' (section 28),
which just means that up to unique isomorphism, there is indeed but \uline{one}

\margin{29}
